% TOP_PROBLEM: {"problemId": "problem.l-p-problem", "canonicalTitle": "L = P Problem", "releaseRank": 60, "primaryDomain": "theoretical_computer_science", "primaryDomainLabel": "Theoretical computer science", "displayStatus": "open", "releaseStatus": "open"}
% !TeX program = lualatex
% Definition and sources only; this file is not an LLM solution attempt.
\documentclass[11pt]{article}
\usepackage[margin=25mm]{geometry}
\usepackage{fontspec}
\setmainfont{Latin Modern Roman}
\usepackage{amsmath,amssymb,mathtools}
\usepackage{unicode-math}
\setmathfont{Latin Modern Math}
\usepackage{xurl,hyperref}
\hypersetup{colorlinks=true,urlcolor=blue}
\setcounter{secnumdepth}{0}
\setlength{\parindent}{0pt}
\setlength{\parskip}{5pt}
\setlength{\emergencystretch}{3em}
\begin{document}
\section*{\texorpdfstring{L = P Problem}{L = P Problem}}\label{60}
\subsection{Definitions and mathematical statement}
In the standard Turing-machine model, $\mathsf L=\mathsf{DSPACE}(\log n)$ counts work-tape space but not read-only input cells. The class $\mathsf P$ consists of total deterministic polynomial-time deciders. The equality question is
\[
 \forall L\in\mathsf P\ \exists\text{ deterministic decider for }L
 \text{ using }O(\log n)\text{ work space}.
\]
The algorithm and constants may depend on the language. A fixed polynomial saving in space, or a logarithmic-space procedure with advice, does not establish the stated equality.
\par\smallskip\textit{Formulation and concept references:} \href{https://www.cs.cornell.edu/courses/cs6820/2022fa/Handouts/CVP.pdf}{[S1]}.

\subsection{Short English statement}
Can every polynomial-time decision problem be solved using only logarithmically much working memory?
\subsection{Sources}
\begin{itemize}
\item[S1]Cornell CS 6820 handout, The Circuit Value Problem, Fall 2022.
\url{https://www.cs.cornell.edu/courses/cs6820/2022fa/Handouts/CVP.pdf}.
\textit{Formulation source.}
\item[S2]Composing Low-Space Algorithms, revision 1 (29 March 2026).
\url{https://eccc.weizmann.ac.il/report/2025/140/revision/1/download/}.
\textit{Further reference; not a proof endorsement.}
\item[S3]Simulating Time With Square-Root Space.
\url{https://arxiv.org/abs/2502.17779}.
\textit{Further reference; not a proof endorsement.}
\end{itemize}
\end{document}
